% ch3.tex -- Module 8, Chapter 3: The quantum atom (Bohr, de Broglie, Schrödinger)
% Syllabus: M8.IQ3.01 - M8.IQ3.05 (NSW Physics Stage 6 Syllabus 2017).
% Figures: figures/ch3_spectra.py, figures/ch3_clouds.py; all other diagrams are TikZ below.
\chapter[The Quantum Atom]{The Quantum Atom: Bohr, de~Broglie and Schrödinger}\label{ch3:chap}
\covers{M8.IQ3.01, M8.IQ3.02, M8.IQ3.03, M8.IQ3.04, M8.IQ3.05}

\inquiry{How is it known that classical physics cannot explain the properties of the atom?}

\begin{goals}
  \item explain why Rutherford's atom fails, and outline Bohr's postulates
  \item use $E = hf$, $E = \dfrac{hc}{\lambda}$ and the Rydberg equation for the lines of
        hydrogen, including the Balmer series
  \item use $\lambda = \dfrac{h}{mv}$ and describe the electron-diffraction evidence for it
  \item assess the limitations of the Rutherford and Bohr models, and analyse what
        Schrödinger added
\end{goals}

%% ==================================================================
\section{From Rutherford's atom to Bohr's energy levels}\label{ch3:sec:bohr}
\covers{M8.IQ3.01, M8.IQ3.03}

\begin{bigidea}
Classical physics says the electron in Rutherford's atom should give out light and crash into
the nucleus. Bohr fixed this with quantum rules. The electron can only have certain energies.
It gives out or takes in light only when it jumps between them.
\end{bigidea}

\begin{keywords}[breakable]
  \kw{Nuclear model}{Rutherford's atom: a tiny, dense, positive nucleus with electrons
    moving around it.}
  \kw{Classical physics}{physics from before quantum ideas: Newton's laws and Maxwell's
    electromagnetism.}
  \kw{Quantised}{only certain values are allowed.}
  \kw{Stationary state}{an allowed orbit. In it the electron does not radiate (give out EM
    waves), so its energy stays the same.}
  \kw{Energy level}{the fixed energy of one stationary state. The lowest is the
    \emph{ground state} ($n = 1$); the others are \emph{excited states}.}
  \kw{Angular momentum}{the amount of turning motion. For an electron moving in a circle it
    is $mvr$.}
\end{keywords}

In \cref{ch2:sec:nucleus} you met the \term{nuclear model}. It is Rutherford's 1911 atom: a tiny,
dense, positive nucleus with electrons moving around it. It explained the gold-foil results.
But \term{classical physics} (Newton's laws and Maxwell's electromagnetism) says it cannot
last.

\begin{refresher}{Ideas you need (Modules 5 and 7)}
\begin{itemize}
  \item An object moving in a circle is accelerating, even at a constant speed.
  \item An \emph{accelerating} charge gives out electromagnetic (EM) waves, which carry energy
        away (Maxwell).
  \item Light comes in packets called \term{photons}. One photon of frequency $f$ has energy
        $E = hf$. Small energies are measured in \term{electronvolts} (eV). One electronvolt
        is the energy an electron gains when it is accelerated through \qty{1}{V}:
        $\qty{1}{eV} = \qty{1.602e-19}{J}$.
\end{itemize}
\end{refresher}

\subsection{What is wrong with Rutherford's atom?}

\begin{sidebyside}[0.38]
  \centering
  \begin{tikzpicture}[x=1cm, y=1cm]
    \draw[hscvel, line width=1pt, ->] plot[domain=0:1000, samples=220, smooth]
      ({1.6*exp(-0.0014*\x)*cos(\x)}, {1.6*exp(-0.0014*\x)*sin(\x)});
    \fill[hscforce] (0,0) circle (3.2pt);
    \node[font=\scriptsize\bfseries, text=white] at (0,0) {$+$};
    \fill[hscvel] (1.6,0) circle (3pt);
    \node[diagramlabel, right] at (1.65,0) {$e^-$};
    \foreach \a/\r in {35/1.36, 150/1.1, 260/0.87} {
      \draw[hscray, line width=0.9pt, ->, decorate,
        decoration={snake, amplitude=1.3pt, segment length=5pt, post length=2.2mm}]
        (\a:\r+0.12) -- (\a:\r+0.95);
    }
    \node[diagramlabel, text=hscray, align=center] at (62:2.45) {EM waves carry\\energy away};
    \node[diagramlabel, below, align=center] at (0,-1.75) {electron spirals in\\(in about
      \qty{e-11}{s})};
  \end{tikzpicture}
  \captionof{figure}{The ``death spiral'' that classical physics predicts.}
  \label{ch3:fig:spiral}
\nextside
  \textbf{Problem 1: it should collapse} (\cref{ch3:fig:spiral}).
  \begin{enumerate}
    \item The orbiting electron is accelerating.
    \item So it gives out EM waves, which carry energy away.
    \item It falls to smaller and smaller orbits. (It actually speeds up as it falls in.
          Its kinetic energy goes up, but its electrical potential energy goes down by twice
          as much. So its total energy falls.)
    \item It should hit the nucleus in about \qty{e-11}{s}.
  \end{enumerate}
  But real atoms last for billions of years. They are \term{stable}: they do not collapse.
\end{sidebyside}

\textbf{Problem 2: the wrong spectrum.} Classically, the light's frequency equals the number of
turns the electron makes each second. As the electron spirals in, it goes round faster and
faster. So the frequency would rise smoothly through every value. That gives a
\term{continuous spectrum} (every colour).

Real hydrogen gives a \term{line spectrum}: only a few sharp wavelengths
(\cref{ch3:fig:spectra}).

Cooler hydrogen gas in front of a hotter light source also \emph{absorbs} light. Shine white
light through it, and dark lines appear. They are at exactly the wavelengths that hot hydrogen
\emph{emits}. (The cooler outer layers of a star do this to the light from the hotter layers
below them. See \cref{ch1:sec:spectra}.)

\begin{figure}
  \centering
  \includegraphics{ch3_spectra.pdf}
  \caption{A continuous spectrum, and hydrogen's emission and absorption spectra.}
  \label{ch3:fig:spectra}
\end{figure}

\textbf{Other problems.} The model does not say what the nucleus is made of, or what holds its
positive charges together. It does not say how the electrons are arranged, so it cannot explain
chemistry.

\subsection{Bohr's postulates (1913)}

Niels Bohr kept the nucleus but added quantum rules. A \term{postulate} is a rule that is
assumed, without proof, to build a model. Different textbooks list Bohr's postulates in
different orders. So learn them by what they say, not by their number:
\begin{enumerate}
  \item Electrons move in circular orbits around the nucleus. The electrostatic (Coulomb)
        attraction holds them there. It is the pull between the positive nucleus and the
        negative electron.
  \item Only certain orbits are allowed. In these \term{stationary states} the electron does
        \emph{not} radiate (give out EM waves), so its energy stays constant.
  \item An electron moves between stationary states only by absorbing or emitting \emph{one}
        photon. The photon's energy equals the energy difference between the two levels:
        $hf = E_\mathrm{i} - E_\mathrm{f}$ for emission, and $hf = E_\mathrm{f} -
        E_\mathrm{i}$ for absorption. ($E_\mathrm{i}$ is the energy of the starting level.
        $E_\mathrm{f}$ is the energy of the final level.)
  \item The electron's \term{angular momentum} is \term{quantised} (only certain values are
        allowed). Angular momentum is the amount of turning motion an orbiting object has.
        For an electron moving in a circle it equals $mvr$ (mass $\times$ speed $\times$
        radius). Bohr's rule is $mvr = \dfrac{nh}{2\pi}$, with $n = 1, 2, 3, \ldots$ So
        $mvr$ can only be a whole-number multiple of $\dfrac{h}{2\pi} = \qty{1.055e-34}{J.s}$.
        (Not on the formula sheet. You may be asked which values are allowed, but not for
        long calculations.)
\end{enumerate}

Think of a ladder. You can stand on any rung, but never between two rungs. Each rung is an
\term{energy level}, named by the number $n$. \Cref{ch3:fig:levels} shows hydrogen's levels.
A jump between two levels is called a \term{transition}. Exam questions often use this word.

\begin{figure}
  \centering
  \begin{tikzpicture}[x=0.9cm, y=0.8cm, font=\footnotesize]
    \foreach \y/\n/\E in {0/1/-13.6, 3.0/2/-3.40, 4.3/3/-1.51, 5.1/4/-0.85, 5.7/5/-0.54,
                          6.2/6/-0.38} {
      \draw[hscink, line width=0.8pt] (0,\y) -- (11.2,\y);
      \node[anchor=west] at (11.3,\y) {$n = \n$ \quad $\num{\E}$ eV};
    }
    \draw[hscink, line width=0.8pt, dashed] (0,6.8) -- (11.2,6.8);
    \node[anchor=west] at (11.3,6.8) {$n = \infty$ \quad 0 eV};
    \node[anchor=east, align=right] at (-0.1,0) {ground\\state};
    \node[anchor=east, align=right] at (-0.1,4.6) {excited\\states};
    \foreach \x/\y in {0.5/3.0, 0.9/4.3, 1.3/5.1, 1.7/5.7} {
      \draw[->, line width=0.9pt, hscfield!80] (\x,\y) -- (\x,0.04);
    }
    \node[draw=hscfield!80, rounded corners=1pt, align=center, fill=white, inner sep=2pt]
      at (3.7,0.65) {Lyman series\\(to $n=1$): ultraviolet};
    \foreach \x/\y/\col/\lab in {4.7/4.3/red!85!black/656, 5.5/5.1/cyan!75!blue/486,
                                 6.3/5.7/blue!80!black/434, 7.1/6.2/violet/410} {
      \draw[->, line width=1.1pt, \col] (\x,\y) -- (\x,3.04);
      \node[below, inner sep=1.5pt, text=\col] at (\x,3.0) {\lab};
    }
    \node[draw=hscink!60, rounded corners=1pt, align=center, fill=white, inner sep=2pt]
      at (6.1,1.95) {Balmer series (to $n=2$): 4 visible\\lines (in nm) + UV lines};
    \foreach \x/\y in {8.4/5.1, 8.8/5.7, 9.2/6.2} {
      \draw[->, line width=0.9pt, red!45!black] (\x,\y) -- (\x,4.34);
    }
    \node[draw=red!45!black, rounded corners=1pt, align=center, fill=white, inner sep=2pt]
      at (8.8,3.6) {Paschen series\\(to $n=3$): infrared};
    \draw[->, line width=0.9pt, hscink, dashed] (10.5,0.04) -- (10.5,6.75);
    \node[align=center, fill=white, inner sep=1pt] at (10.5,1.6)
      {ionisation\\from $n=1$:\\13.6 eV};
  \end{tikzpicture}
  \caption{Energy levels of hydrogen (spacing not to scale) and the three main series of
    lines (\cref{ch3:sec:rydberg}).}
  \label{ch3:fig:levels}
\end{figure}

\subsubsection{Why are the energies negative?}
We choose 0 eV for an electron that is just free ($n = \infty$). An electron held by the
nucleus (a \emph{bound} electron) has less energy than a free one. You must add energy to free
it. So its energy is below zero: negative. (Each level is $E_n = -\dfrac{\qty{13.6}{eV}}{n^2}$.
This is not on the formula sheet. Questions usually give the levels.)

\subsubsection{Conservation of energy links the levels to the light}
\begin{sidebyside}[0.4]
  \centering
  \begin{tikzpicture}[x=1cm, y=1cm, font=\small]
    \foreach \yy/\title in {2.45/Emission, 0/Absorption} {
      \begin{scope}[shift={(0,\yy)}]
        \draw[hscink, line width=0.8pt] (0,0) -- (1.8,0) (0,1.2) -- (1.8,1.2);
        \node[anchor=east, font=\footnotesize] at (-0.08,1.2) {upper};
        \node[anchor=east, font=\footnotesize] at (-0.08,0) {lower};
        \node[font=\small\bfseries] at (0.9,1.6) {\title};
      \end{scope}
    }
    % emission: electron drops, photon leaves
    \fill[hscvel] (0.45,3.65) circle (2.4pt);
    \draw[->, hscvel, line width=1.1pt] (0.45,3.57) -- (0.45,2.53);
    \draw[hscray, line width=0.9pt, ->, decorate,
      decoration={snake, amplitude=1.6pt, segment length=6pt, post length=2.2mm}]
      (0.7,3.05) -- (2.5,3.05);
    \node[anchor=west, text=hscray, align=left] at (2.52,3.05) {photon\\out};
    % absorption: photon arrives, electron jumps up
    \draw[hscray, line width=0.9pt, ->, decorate,
      decoration={snake, amplitude=1.6pt, segment length=6pt, post length=2.2mm}]
      (2.5,0.6) -- (0.7,0.6);
    \node[anchor=west, text=hscray, align=left] at (2.52,0.6) {photon\\in};
    \fill[hscvel] (0.45,0) circle (2.4pt);
    \draw[->, hscvel, line width=1.1pt] (0.45,0.08) -- (0.45,1.12);
  \end{tikzpicture}
  \captionof{figure}{In both, the photon's energy equals the gap between the two levels.}
  \label{ch3:fig:emitabsorb}
\nextside
  \begin{itemize}
    \item \textbf{Emission:} an electron drops from $E_\mathrm{i}$ to a lower level
          $E_\mathrm{f}$. One photon carries away the energy lost: $hf = E_\mathrm{i} -
          E_\mathrm{f}$.
    \item \textbf{Absorption:} a photon is absorbed only if its energy matches a gap
          exactly. The electron jumps up. Other photons with less than the ionisation
          energy pass straight through.
    \item \textbf{Ionisation:} the electron gains enough energy to reach $n = \infty$ and
          leave. The smallest energy needed is the \term{ionisation energy} (\qty{13.6}{eV}
          from $n = 1$). Any extra becomes the free electron's kinetic energy.
  \end{itemize}
  Only certain gaps exist, so only certain photon energies and frequencies come out
  (\cref{ch3:fig:emitabsorb}). This is why hydrogen gives a line spectrum.
\end{sidebyside}

\begin{toolkit}{Changing between eV and J}
eV $\to$ J: multiply by \num{1.602e-19}. \quad J $\to$ eV: divide by \num{1.602e-19}.\\
Example: $\qty{10.2}{eV} \times \num{1.602e-19} = \qty{1.63e-18}{J}$.
\end{toolkit}

\begin{formulacard}{Photon energy}
  \cardformula{E = hf}
  \begin{symbols}
    \sym{$E$}{energy of one photon}{\si{J}}
    \sym{$h$}{Planck constant, \num{6.626e-34}}{\si{J.s}}
    \sym{$f$}{frequency of the light}{\si{Hz} ($\si{s^{-1}}$)}
  \end{symbols}
  \inwords{double the frequency and each photon has double the energy.}
  \whentouse{a photon's energy from its frequency, or the photon's frequency from an energy
    gap.}
  \whennottouse{for the energy of one level, or with $E$ still in eV.}
  \rearranged{$f = \dfrac{E}{h}$ \qquad $hf = E_\mathrm{i} - E_\mathrm{f}$ (emission; for
    absorption use $E_\mathrm{f} - E_\mathrm{i}$)}
  \tensecond{$f = \qty{5.0e14}{Hz}$, so
    $E = \num{6.626e-34} \times \num{5.0e14} = \qty{3.3e-19}{J}$ (\qty{2.1}{eV}).}
  \trap{the photon energy is the \emph{gap} between two levels, not one level's energy.
    Change eV to J before you divide by $h$.}
\end{formulacard}

\begin{workedexample}{A photon from a big drop}\label{ch3:we:fiveone}
Use \cref{ch3:fig:levels}. An electron in hydrogen drops from $n = 5$ to $n = 1$. Calculate
the photon's energy (in eV and J), its frequency and its wavelength.
\given{$E_5 = \qty{-0.54}{eV}$, \ $E_1 = \qty{-13.6}{eV}$, \ $h = \qty{6.626e-34}{J.s}$,
  \ $c = \qty{3.00e8}{m.s^{-1}}$}
\find{the photon energy $E$, its frequency $f$ and its wavelength $\lambda$}
\formula{$E = E_\mathrm{i} - E_\mathrm{f}$, \ then $f = \dfrac{E}{h}$ (from $E = hf$), \ then
  $\lambda = \dfrac{c}{f}$ (from $c = f\lambda$)}
\substitute{\eqlines{E &= \qty{-0.54}{eV} - (\qty{-13.6}{eV}) = \qty{13.06}{eV}\\
  E &= 13.06 \times \qty{1.602e-19}{J} = \qty{2.092e-18}{J} \quad\text{(eV to J)}\\
  f &= \frac{\qty{2.092e-18}{J}}{\qty{6.626e-34}{J.s}} = \qty{3.157e15}{Hz}\\
  \lambda &= \frac{\qty{3.00e8}{m.s^{-1}}}{\qty{3.157e15}{Hz}} = \qty{9.50e-8}{m}
    = \qty{95.0}{nm}}}
\answer{\final{$E = \qty{13.1}{eV} = \qty{2.09e-18}{J}$}, \ \final{$f = \qty{3.16e15}{Hz}$}
  and \final{$\lambda = \qty{95.0}{nm}$}}
\checkstep{a drop to $n = 1$ is a Lyman line, so it should be ultraviolet. Visible light is
  about \qtyrange{4.3e14}{7.5e14}{Hz} (\qtyrange{700}{400}{nm}). Our $f$ is higher and our
  $\lambda$ is shorter \cmark.}
\end{workedexample}

\begin{workedexample}{Can red light ionise an excited atom?}\label{ch3:we:ionise}
A hydrogen atom is in the $n = 3$ level ($E_3 = \qty{-1.51}{eV}$). Red light of frequency
\qty{4.57e14}{Hz} shines on it. Can one photon ionise the atom? If so, find the kinetic
energy of the freed electron.
\given{$E_3 = \qty{-1.51}{eV}$, \ $f = \qty{4.57e14}{Hz}$, \ $h = \qty{6.626e-34}{J.s}$}
\find{the photon energy; compare it with the ionisation energy from $n = 3$}
\formula{$E = hf$; \ ionisation energy $E_\mathrm{ion} = 0 - E_3$ (a free electron, at
  $n = \infty$, has \qty{0}{eV}); \ kinetic energy $K = E - E_\mathrm{ion}$}
\substitute{\eqlines{E &= \qty{6.626e-34}{J.s} \times \qty{4.57e14}{Hz} = \qty{3.028e-19}{J}\\
  E &= \frac{\qty{3.028e-19}{J}}{\num{1.602e-19}} = \qty{1.89}{eV} \quad\text{(J to eV)}\\
  E_\mathrm{ion} &= 0 - (\qty{-1.51}{eV}) = \qty{1.51}{eV}\\
  K &= \qty{1.89}{eV} - \qty{1.51}{eV} = \qty{0.38}{eV}\\
  K &= 0.38 \times \qty{1.602e-19}{J} = \qty{6.1e-20}{J} \quad\text{(eV to J)}}}
\answer{yes, because $\qty{1.89}{eV} > \qty{1.51}{eV}$.
  \final{$K = \qty{0.38}{eV} = \qty{6.1e-20}{J}$}}
\checkstep{energy is conserved: $\qty{1.51}{eV} + \qty{0.38}{eV} = \qty{1.89}{eV}$ \cmark.
  The same photon could not ionise a ground-state atom (that needs \qty{13.6}{eV}).}
\end{workedexample}

\begin{mistakes}
\begin{itemize}
  \item \xmark\ ``Rutherford's electron slows down as it spirals in.'' \cmark\ It speeds up;
        its \emph{total} energy falls.
  \item \xmark\ Using one level's energy as the photon energy. \cmark\ Use the difference,
        $E_\mathrm{i} - E_\mathrm{f}$, and change eV to J before using $E = hf$.
  \item \xmark\ ``An atom absorbs any photon with enough energy.'' \cmark\ A jump needs an
        exact match to a gap. Only ionisation takes any larger energy.
\end{itemize}
\end{mistakes}

\begin{examtip}
``Explain why Rutherford's model is unstable'' (3 marks). Write three linked points:
\begin{enumerate}
  \item The orbiting electron is accelerating.
  \item Accelerating charges give out EM waves (Maxwell), so the electron loses energy.
  \item It should spiral into the nucleus, but real atoms are stable.
\end{enumerate}
``Assess the limitations of Rutherford's model'' (4--5 marks). First give one success: the
gold-foil results show a tiny, dense, positive nucleus (still accepted today). Then give
limitations, each linked to evidence:
\begin{enumerate}
  \item Classical physics predicts collapse in about \qty{e-11}{s}, but atoms are stable.
  \item It predicts a continuous spectrum, but hydrogen gives a line spectrum.
  \item It does not say what the nucleus is made of, or how the electrons are arranged.
\end{enumerate}
End with a judgement, e.g. \emph{right about the nucleus, but it cannot explain stable atoms
or line spectra, so quantum ideas (Bohr) were needed.}

``Outline Bohr's postulates'' (4 marks): one mark for each, by content. Always include
``electrons in stationary states do not radiate''.
\end{examtip}

\begin{quickcheck}
  \item Explain why classical physics predicts that Rutherford's atom would collapse, and why
        it predicts a continuous spectrum.
  \item Use \cref{ch3:fig:levels} to find the energy of the photon emitted when an electron
        drops from $n = 6$ to $n = 3$. Give it in eV and in J.
  \item Hydrogen atoms are excited to $n = 4$. (a) How many different spectral lines can
        they emit? (b) Which jump gives the highest frequency?
  \item A photon of energy \qty{11.0}{eV} hits a hydrogen atom in its ground state. Explain
        whether the photon can be absorbed.
\end{quickcheck}

%% ==================================================================
\section{Hydrogen's line spectrum and the Rydberg equation}\label{ch3:sec:rydberg}
\covers{M8.IQ3.02, M8.IQ3.03}

\begin{bigidea}
Each line in hydrogen's spectrum comes from one kind of jump between two energy levels. The
Rydberg equation gives the wavelength of every line. Jumps that end on $n = 2$ form the Balmer
series; four of these lines are visible.
\end{bigidea}

\begin{keywords}
  \kw{Balmer series}{the hydrogen lines made by jumps that end on $n = 2$.}
  \kw{Lyman / Paschen series}{the lines made by jumps that end on $n = 1$ (ultraviolet) /
    on $n = 3$ (infrared).}
  \kw{Rydberg constant ($R$)}{a fixed number for hydrogen, \qty{1.097e7}{m^{-1}}. It sets
    the size of all of hydrogen's wavelengths.}
\end{keywords}

Each element gives its own set of lines, like a fingerprint. This is how astronomers can tell
that a star contains hydrogen.

In 1885 Johann Balmer found a formula that fitted hydrogen's four visible lines. In 1888
Johannes Rydberg made it general. Neither knew \emph{why} it worked. Bohr's model gave the
reason: each line is one jump, and the photon carries the energy difference. Bohr even
calculated $R$ from $h$ and the electron's mass and charge.

\Cref{ch3:tab:series} lists the three main series (see also \cref{ch3:fig:levels}). Each
series has a \term{series limit}: its shortest wavelength, for a jump that starts at
$n = \infty$.

\begin{table}
  \caption{The three main series of hydrogen lines.}\label{ch3:tab:series}
  \centering
  \begin{tabular}{@{}L{20mm}L{18mm}L{30mm}L{66mm}@{}}
    \toprule
    \colhead{Series} & \colhead{Ends on} & \colhead{Region} & \colhead{Wavelengths}\\
    \midrule
    Lyman   & $n_\mathrm{f} = 1$ & ultraviolet & \qtyrange{91}{122}{nm}\\
    Balmer  & $n_\mathrm{f} = 2$ & visible, near UV & 656, 486, 434, \qty{410}{nm} (visible),
              then UV lines down to \qty{365}{nm}\\
    Paschen & $n_\mathrm{f} = 3$ & infrared & \qtyrange{820}{1875}{nm}\\
    \bottomrule
  \end{tabular}
\end{table}

Regions of the EM spectrum used in this chapter: \textbf{ultraviolet} (UV) is shorter than
about \qty{400}{nm}. \textbf{Visible} light is about \qty{400}{nm} (violet) to \qty{700}{nm}
(red). \textbf{Infrared} (IR) is longer than about \qty{700}{nm}.

\begin{formulacard}[combine={$E = hf$ and $v = f\lambda$}]{Photon energy from wavelength}
  \cardformula{E = \frac{hc}{\lambda}}
  \begin{symbols}
    \sym{$E$}{energy of one photon}{\si{J}}
    \sym{$h$}{Planck constant, \num{6.626e-34}}{\si{J.s}}
    \sym{$c$}{speed of light, \num{3.00e8}}{\si{m.s^{-1}}}
    \sym{$\lambda$}{wavelength of the light}{\si{m}}
  \end{symbols}
  \inwords{the shorter the wavelength, the more energy each photon has.}
  \whentouse{to change between a photon's wavelength and its energy.}
  \whennottouse{for electrons or other particles with mass (use $\lambda = \dfrac{h}{mv}$).}
  \rearranged{$\lambda = \dfrac{hc}{E}$ \qquad $E = hc \times \dfrac{1}{\lambda}$ (if you
    already have $\dfrac{1}{\lambda}$)}
  \tensecond{$\lambda = \qty{500}{nm} = \qty{5.00e-7}{m}$, so
    $E = \dfrac{\num{6.626e-34} \times \num{3.00e8}}{\num{5.00e-7}} = \qty{3.98e-19}{J}$.}
  \derive{two formulas from the sheet. For light, the wave speed $v$ is $c$.
    \begin{align*}
      c &= f\lambda \ \Rightarrow\ f = \frac{c}{\lambda}
        && \text{$v = f\lambda$ for light; divide by $\lambda$}\\
      E &= hf = \frac{hc}{\lambda} && \text{replace $f$ in $E = hf$}
    \end{align*}}
  \trap{change nm to m ($\times 10^{-9}$). If you have $\dfrac{1}{\lambda}$, \emph{multiply}
    it by $hc$.}
\end{formulacard}

\textbf{Where the Rydberg equation comes from.} Join two ideas. Energy conservation says the
photon's energy is the gap: $\dfrac{hc}{\lambda} = E_\mathrm{i} - E_\mathrm{f}$. Hydrogen's
levels are $E_n = -\dfrac{\qty{13.6}{eV}}{n^2}$ (not on the formula sheet).
\begin{align*}
  \frac{hc}{\lambda} &= \left(-\frac{\qty{13.6}{eV}}{n_\mathrm{i}^2}\right)
    - \left(-\frac{\qty{13.6}{eV}}{n_\mathrm{f}^2}\right)
    && \text{put the two levels into $E_\mathrm{i} - E_\mathrm{f}$}\\
  \frac{hc}{\lambda} &= \qty{13.6}{eV} \times \left(\frac{1}{n_\mathrm{f}^2}
    - \frac{1}{n_\mathrm{i}^2}\right) && \text{take out \qty{13.6}{eV}; the minus signs swap
    the order}\\
  \frac{1}{\lambda} &= \frac{\qty{13.6}{eV}}{hc} \times \left(\frac{1}{n_\mathrm{f}^2}
    - \frac{1}{n_\mathrm{i}^2}\right) && \text{divide both sides by $hc$}
\end{align*}
This is the Rydberg equation, with $R = \dfrac{\qty{13.6}{eV}}{hc}$.

Check the numbers: $hcR = \num{6.626e-34} \times \num{3.00e8} \times \num{1.097e7} =
\qty{2.18e-18}{J}$, and $\dfrac{\qty{2.18e-18}{J}}{\num{1.602e-19}} = \qty{13.6}{eV}$ \cmark.

Test it on \cref{ch3:we:fiveone} ($5 \to 1$): $\dfrac{1}{\lambda} = \num{1.097e7} \times
\left(\dfrac{1}{1^2} - \dfrac{1}{5^2}\right) = \qty{1.053e7}{m^{-1}}$, so $\lambda =
\qty{95.0}{nm}$. This is the same answer we found from the energy levels \cmark.

\begin{toolkit}{Fractions with squares, and reciprocals}
Square first, then subtract: $\dfrac{1}{3^2} - \dfrac{1}{6^2} = \dfrac19 - \dfrac{1}{36} =
0.1111 - 0.0278 = 0.0833$. The equation gives $\dfrac{1}{\lambda}$; press $x^{-1}$ to get
$\lambda$. Example: $\dfrac{1}{\lambda} = \qty{2.0e6}{m^{-1}}$ gives $\lambda =
\qty{5.0e-7}{m} = \qty{500}{nm}$ (because $\qty{1}{nm} = \qty{e-9}{m}$).
\end{toolkit}

\begin{formulacard}{The Rydberg equation (hydrogen)}
  \cardformula{\frac{1}{\lambda} = R\left(\frac{1}{n_\mathrm{f}^2} - \frac{1}{n_\mathrm{i}^2}\right)}
  \begin{symbols}
    \sym{$\lambda$}{wavelength of the emitted (or absorbed) photon}{\si{m}}
    \sym{$R$}{Rydberg constant for hydrogen, \num{1.097e7}}{\si{m^{-1}}}
    \sym{$n_\mathrm{f}$}{number of the final (end) level}{no unit}
    \sym{$n_\mathrm{i}$}{number of the initial (start) level}{no unit}
  \end{symbols}
  \inwords{the wavelength of a hydrogen line depends only on where the jump starts and ends.}
  \whentouse{any hydrogen line. Emission: $n_\mathrm{i} > n_\mathrm{f}$. For a series limit,
    put $n_\mathrm{i} = \infty$, so $\dfrac{1}{n_\mathrm{i}^2} = 0$. For ionisation, the free
    electron's level is the $\infty$ one.}
  \whennottouse{atoms with more than one electron (helium, sodium, \ldots).}
  \rearranged{$\lambda = 1 \div \left[R\left(\dfrac{1}{n_\mathrm{f}^2} -
    \dfrac{1}{n_\mathrm{i}^2}\right)\right]$ \qquad
    $E = hcR\left(\dfrac{1}{n_\mathrm{f}^2} - \dfrac{1}{n_\mathrm{i}^2}\right)$}
  \tensecond{$6 \to 3$: $\dfrac{1}{\lambda} = \num{1.097e7} \times \left(\dfrac19 -
    \dfrac{1}{36}\right) = \qty{9.14e5}{m^{-1}}$, so $\lambda = \qty{1094}{nm}$ (infrared).}
  \trap{square the $n$ values. For absorption the bracket comes out negative. Ignore the minus
    sign: the wavelength is the same as for the reverse drop. Here $R$ is not resistance.}
\end{formulacard}

\begin{workedexample}{The violet Balmer line}\label{ch3:we:sixtwo}
An electron in hydrogen drops from $n = 6$ to $n = 2$. Calculate the wavelength of the photon
and state its colour. Then find the photon's energy in J and eV.
\given{$n_\mathrm{i} = 6$, \ $n_\mathrm{f} = 2$, \ $R = \qty{1.097e7}{m^{-1}}$,
  \ $h = \qty{6.626e-34}{J.s}$, \ $c = \qty{3.00e8}{m.s^{-1}}$}
\find{the wavelength $\lambda$ and the photon energy $E$}
\formula{$\dfrac{1}{\lambda} = R\left(\dfrac{1}{n_\mathrm{f}^2} -
  \dfrac{1}{n_\mathrm{i}^2}\right)$ \quad then \quad $E = \dfrac{hc}{\lambda}$}
\substitute{\eqlines{\frac{1}{\lambda} &= \qty{1.097e7}{m^{-1}} \times
    \left(\frac{1}{4} - \frac{1}{36}\right)\\
  &= \qty{1.097e7}{m^{-1}} \times 0.2222 = \qty{2.438e6}{m^{-1}}\\
  \lambda &= \frac{1}{\qty{2.438e6}{m^{-1}}} = \qty{4.102e-7}{m} = \qty{410.2}{nm}\\
  E &= \frac{\qty{6.626e-34}{J.s} \times \qty{3.00e8}{m.s^{-1}}}{\qty{4.102e-7}{m}}
    = \qty{4.846e-19}{J}\\
  E &= \frac{\qty{4.846e-19}{J}}{\num{1.602e-19}} = \qty{3.02}{eV} \quad\text{(J to eV)}}}
\answer{\final{$\lambda = \qty{410}{nm}$}, violet. \ \final{$E = \qty{4.85e-19}{J} =
  \qty{3.02}{eV}$}}
\checkstep{it ends on $n = 2$, so it is a Balmer line, and \qty{410}{nm} is one of the four
  visible lines \cmark. Violet photons carry more energy than red ones (\qty{1.89}{eV}).}
\end{workedexample}

\begin{workedexample}{Which jump made this line?}\label{ch3:we:paschen}
Hydrogen gives an infrared line at \qty{1282}{nm}. (a) Explain why it cannot be a Balmer
line. (b) It ends on $n = 3$. Find the level where the jump started.
\given{$\lambda = \qty{1282}{nm} = \qty{1.282e-6}{m}$, \ $R = \qty{1.097e7}{m^{-1}}$,
  \ $n_\mathrm{f} = 3$}
\find{(a) a reason; (b) $n_\mathrm{i}$}
\formula{$\dfrac{1}{R\lambda} = \dfrac{1}{n_\mathrm{f}^2} - \dfrac{1}{n_\mathrm{i}^2}$
  (divide the Rydberg equation by $R$). Add $\dfrac{1}{n_\mathrm{i}^2}$ to both sides, then
  take $\dfrac{1}{R\lambda}$ from both sides:
  $\dfrac{1}{n_\mathrm{i}^2} = \dfrac{1}{n_\mathrm{f}^2} - \dfrac{1}{R\lambda}$}
\substitute{\eqlines{\frac{1}{R\lambda} &= \frac{1}{\qty{1.097e7}{m^{-1}} \times
    \qty{1.282e-6}{m}} = 0.07111 \quad\text{(no unit)}\\
  \frac{1}{n_\mathrm{i}^2} &= \frac{1}{9} - 0.07111 = 0.11111 - 0.07111 = 0.04000\\
  n_\mathrm{i}^2 &= \frac{1}{0.04000} = 25.0 \quad\text{(take the reciprocal)}\\
  n_\mathrm{i} &= \sqrt{25.0} = 5 \quad\text{(square root)}}}
\answer{(a) The longest Balmer wavelength comes from the smallest Balmer jump, $3 \to 2$. A
  smaller jump gives a smaller $\dfrac{1}{\lambda}$, so a longer $\lambda$. That longest line is
  \qty{656}{nm} (red, \cref{ch3:tab:series}). So every Balmer line is \qty{656}{nm} or
  shorter. \qty{1282}{nm} is longer than \qty{656}{nm}, so it cannot be a Balmer line.\newline
  (b) \final{the jump is $n = 5 \to n = 3$} (Paschen series).}
\checkstep{$n_\mathrm{i}$ is a whole number, as it must be for a real hydrogen line \cmark.}
\end{workedexample}

\begin{mistakes}
\begin{itemize}
  \item \xmark\ Forgetting to square: $\frac12 - \frac16$. \cmark\ $\frac{1}{2^2} -
        \frac{1}{6^2}$.
  \item \xmark\ Stopping at $\dfrac{1}{\lambda}$. \cmark\ Take the reciprocal to get
        $\lambda$.
  \item \xmark\ ``All Balmer lines are visible.'' \cmark\ Only four are; the rest are UV.
\end{itemize}
\end{mistakes}

\begin{examtip}
``Define the Balmer series'' (1 mark): \emph{the lines produced when electrons in hydrogen
drop to $n = 2$}. To ``relate the energy levels to the spectrum'', use energy conservation.
\emph{The photon carries away exactly the energy the atom loses: $hf = E_\mathrm{i} -
E_\mathrm{f}$. Only certain gaps exist, so only certain wavelengths appear.}
\end{examtip}

\begin{quickcheck}
  \item Calculate the wavelength of the photon emitted when an electron drops from $n = 2$ to
        $n = 1$. In which part of the EM spectrum is it?
  \item Find the shortest wavelength in the Balmer series (a jump from $n = \infty$ to
        $n = 2$). Is it visible?
  \item Ground-state hydrogen atoms absorb light of wavelength \qty{102.6}{nm}. To which level
        do the electrons jump?
\end{quickcheck}

%% ==================================================================
\section{De Broglie's matter waves}\label{ch3:sec:debroglie}
\covers{M8.IQ3.04, M8.IQ3.01}

\begin{bigidea}
Light, a wave, can act like particles (photons). In 1924 Louis de Broglie suggested the
reverse: particles such as electrons can act like waves, with wavelength $\lambda =
\dfrac{h}{mv}$. Electron diffraction experiments confirmed it, and the idea explains Bohr's
orbits.
\end{bigidea}

\begin{keywords}
  \kw{Matter wave}{the wave that goes with a moving particle.}
  \kw{Wave--particle duality}{light and matter both show wave and particle behaviour.}
  \kw{Momentum ($p$)}{mass times velocity, $p = mv$, in \si{kg.m.s^{-1}}.}
\end{keywords}

\begin{refresher}{Diffraction and standing waves (Year 11)}
\textbf{Diffraction:} a wave spreads out after a gap. The spreading is easy to see only when
the gap is about the same size as the wavelength (or smaller). Only waves diffract.
\textbf{Standing wave:} two equal waves travelling in opposite directions make a pattern that
does not travel. It has \emph{nodes}: points that never move. On a string fixed at both ends,
a whole number of half-wavelengths fits: $\lambda = \dfrac{2L}{n}$.
\end{refresher}

\textbf{De Broglie's idea.} Light can push on things. (For example, sunlight pushes on a
spacecraft's solar sail.) So photons carry momentum. A photon of wavelength $\lambda$ has
momentum $p = \dfrac{h}{\lambda}$ (not on the formula sheet). De Broglie guessed that nature
works the same way in both directions. So the same rule should hold for particles, with
$p = mv$. That gives $\lambda = \dfrac{h}{mv}$. This was a \term{hypothesis}, an idea to be
tested.

\textbf{Why we never see a ball diffract.} $h$ is tiny. For a ball, $mv$ is large, so
$\lambda$ is far too small to notice. For an electron at lab speeds (about
\qtyrange{e6}{e7}{m.s^{-1}}), $m$ is tiny, so $\lambda \approx \qty{e-10}{m}$. That is about
the gap between atoms in a crystal. So crystals can diffract electrons. In 1927 Davisson and
Germer, and G.\,P. Thomson, saw this, and the wavelengths matched $\dfrac{h}{mv}$
(\cref{ch3:sec:experiments}).

\textbf{Standing waves explain Bohr's orbits.} Wrap the electron wave around an orbit
(\cref{ch3:fig:orbitwaves}). If a whole number of wavelengths fits, the wave joins itself
smoothly: a \term{standing wave}, and the orbit is allowed. If not, the wave overlaps itself out
of step and cancels, so the orbit cannot exist. A standing wave does not travel, so it does not
carry energy away. This gives a reason for Bohr's ``no radiation'' postulate.

\begin{figure}
  \centering
  \begin{tikzpicture}[x=1cm, y=1cm]
    \foreach \n/\xx in {3/0, 4/3.4, 5/6.8} {
      \begin{scope}[shift={(\xx,0)}]
        \draw[hscmuted!70, dashed, line width=0.5pt] (0,0) circle (0.75);
        \fill[hscforce] (0,0) circle (1.8pt);
        \draw[hscseriesA, line width=1.1pt] plot[domain=0:360, samples=240, smooth]
          ({(0.75+0.14*sin(\n*\x))*cos(\x)}, {(0.75+0.14*sin(\n*\x))*sin(\x)});
        \node[diagramlabel, align=center] at (0,-1.45) {$n = \n$: \n\ whole\\waves fit};
      \end{scope}
    }
    \begin{scope}[shift={(10.2,0)}]
      \draw[hscmuted!70, dashed, line width=0.5pt] (0,0) circle (0.75);
      \fill[hscforce] (0,0) circle (1.8pt);
      \draw[hscmistake, line width=0.8pt] plot[domain=0:720, samples=480, smooth]
        ({(0.75+0.14*sin(3.4*\x))*cos(\x)}, {(0.75+0.14*sin(3.4*\x))*sin(\x)});
      \node[diagramlabel, align=center] at (0,-1.45) {3.4 waves: out\\of step, cancels};
    \end{scope}
  \end{tikzpicture}
  \caption{Electron waves around orbits (not to scale). Blue: three allowed orbits
    ($n = 3$, 4 and 5). Red: not allowed; it is drawn for two laps, and each lap is out of step
    with the one before.}
  \label{ch3:fig:orbitwaves}
\end{figure}

\begin{formulacard}{de Broglie wavelength}
  \cardformula{\lambda = \frac{h}{mv}}
  \begin{symbols}
    \sym{$\lambda$}{de Broglie wavelength of the particle}{\si{m}}
    \sym{$h$}{Planck constant, \num{6.626e-34}}{\si{J.s}}
    \sym{$m$}{mass of the particle}{\si{kg}}
    \sym{$v$}{speed of the particle}{\si{m.s^{-1}}}
  \end{symbols}
  \inwords{the more momentum a particle has, the shorter its wavelength.}
  \whentouse{any moving particle with mass (electron, proton, neutron, ball).}
  \whennottouse{photons (use $E = hc/\lambda$); speeds above about $0.1c$ (use
    $\lambda = h/p_\mathrm{v}$, with $p_\mathrm{v}$ from the sheet).}
  \rearranged{$v = \dfrac{h}{m\lambda}$ \qquad $m = \dfrac{h}{\lambda v}$ \qquad
    $p = mv = \dfrac{h}{\lambda}$}
  \tensecond{a proton ($\qty{1.673e-27}{kg}$) at \qty{1.0e4}{m.s^{-1}}:
    $\lambda = \dfrac{\num{6.626e-34}}{\num{1.673e-27} \times \num{1.0e4}} =
    \qty{4.0e-11}{m}$.}
  \derive{questions may ask you to ``derive'' it from $E = mc^2$ and $E = hf$. This is a
    memory aid, not a proof. A photon has no rest mass, so here $mc$ just stands for the
    photon's momentum $p$ (the middle line says $p = h/\lambda$).
    \begin{align*}
      mc^2 &= hf = \frac{hc}{\lambda} && \text{write the photon's energy as $mc^2$; $f =
        \tfrac{c}{\lambda}$}\\
      \lambda &= \frac{hc}{mc^2} = \frac{h}{mc} && \text{multiply by $\lambda$, divide by
        $mc^2$, cancel one $c$}\\
      \lambda &= \frac{h}{mv} && \text{de Broglie's assumption: same rule for a particle}
    \end{align*}}
  \trap{change \si{km.h^{-1}} to \si{m.s^{-1}} ($\div 3.6$). Faster means a \emph{shorter}
    wavelength. This $\lambda$ is not the decay constant.}
\end{formulacard}

\begin{formulacard}[notonsheet]{Standing electron wave in an orbit}
  \cardformula{2\pi r = n\lambda}
  \begin{symbols}
    \sym{$r$}{radius of the orbit}{\si{m}}
    \sym{$\lambda$}{de Broglie wavelength of the electron}{\si{m}}
    \sym{$n$}{a whole number, the level number (1, 2, 3, \ldots)}{no unit}
  \end{symbols}
  \inwords{an allowed orbit's circumference holds a whole number of electron wavelengths.}
  \whentouse{to explain why only some Bohr orbits are allowed, or to find $\lambda$ or $r$ for
    level $n$.}
  \whennottouse{in Schrödinger's model, which has no orbits.}
  \rearranged{$\lambda = \dfrac{2\pi r}{n}$ \qquad $r = \dfrac{n\lambda}{2\pi}$ \qquad
    $mvr = \dfrac{nh}{2\pi}$}
  \tensecond{$n = 2$, $r = \qty{2.12e-10}{m}$:
    $\lambda = \dfrac{2\pi \times \num{2.12e-10}}{2} = \qty{6.66e-10}{m}$.}
  \derive{put de Broglie's wavelength into the standing-wave rule.
    \begin{align*}
      2\pi r &= n\frac{h}{mv} && \text{replace $\lambda$ with $\dfrac{h}{mv}$}\\
      mvr &= \frac{nh}{2\pi} && \text{multiply both sides by $\dfrac{mv}{2\pi}$}
    \end{align*}
    This is Bohr's angular-momentum postulate. De Broglie's waves \emph{explain} it.}
  \trap{$n$ must be a whole number. The circumference is $2\pi r$ (radius, not diameter).}
\end{formulacard}

\begin{workedexample}{An electron and a cricket ball}\label{ch3:we:ball}
Find the de Broglie wavelength of (a) an electron moving at \qty{4.0e6}{m.s^{-1}}, and (b) a
\qty{0.16}{kg} cricket ball bowled at \qty{40}{m.s^{-1}}. Which one could a crystal diffract
(atoms about \qty{e-10}{m} apart)?
\given{$m_\mathrm{e} = \qty{9.109e-31}{kg}$, $v_\mathrm{e} = \qty{4.0e6}{m.s^{-1}}$; \
  $m_\mathrm{b} = \qty{0.16}{kg}$, $v_\mathrm{b} = \qty{40}{m.s^{-1}}$; \
  $h = \qty{6.626e-34}{J.s}$}
\find{$\lambda$ for each, then compare with \qty{e-10}{m}}
\formula{$\lambda = \dfrac{h}{mv}$}
\substitute{\eqlines{\lambda_\mathrm{e} &= \frac{\qty{6.626e-34}{J.s}}{\qty{9.109e-31}{kg}
    \times \qty{4.0e6}{m.s^{-1}}} = \qty{1.82e-10}{m}\\
  \lambda_\mathrm{b} &= \frac{\qty{6.626e-34}{J.s}}{\qty{0.16}{kg} \times \qty{40}{m.s^{-1}}}
    = \qty{1.04e-34}{m}}}
\answer{(a) \final{$\lambda_\mathrm{e} = \qty{1.8e-10}{m}$} \ (b)
  \final{$\lambda_\mathrm{b} = \qty{1.0e-34}{m}$}. Only the electron's wavelength is about
  the atom spacing, so only the electron can be diffracted.}
\checkstep{units: $\unit{J} = \unit{kg.m^2.s^{-2}}$, so $\unit{J.s} = \unit{kg.m^2.s^{-1}}$.
  Then $\unit{kg.m^2.s^{-1}} \div \unit{kg.m.s^{-1}} = \unit{m}$ \cmark. The ball's $\lambda$
  is far smaller than a nucleus (\qty{e-15}{m}).}
\end{workedexample}

\begin{workedexample}{Waves in Bohr's third orbit}\label{ch3:we:third}
In Bohr's model, hydrogen's $n = 3$ orbit has radius \qty{4.76e-10}{m}, and the electron moves
at \qty{7.29e5}{m.s^{-1}}. Show that exactly three electron wavelengths fit around this orbit.
\given{$r = \qty{4.76e-10}{m}$, \ $v = \qty{7.29e5}{m.s^{-1}}$, \ $m = \qty{9.109e-31}{kg}$,
  \ $h = \qty{6.626e-34}{J.s}$}
\find{the number of wavelengths around the orbit, $\dfrac{2\pi r}{\lambda}$}
\formula{$\lambda = \dfrac{h}{mv}$ \quad then \quad $n = \dfrac{2\pi r}{\lambda}$ (from
  $2\pi r = n\lambda$)}
\substitute{\eqlines{\lambda &= \frac{\qty{6.626e-34}{J.s}}{\qty{9.109e-31}{kg} \times
    \qty{7.29e5}{m.s^{-1}}} = \qty{9.978e-10}{m}\\
  n &= \frac{2\pi \times \qty{4.76e-10}{m}}{\qty{9.978e-10}{m}} =
    \frac{\qty{2.991e-9}{m}}{\qty{9.978e-10}{m}} = 2.998}}
\answer{\final{$n = 3.00$} (3 significant figures): three whole wavelengths fit.}
\checkstep{no unit (\unit{m} $\div$ \unit{m}) \cmark. It equals the level number, as de Broglie
  predicts.}
\end{workedexample}

\begin{mistakes}
\begin{itemize}
  \item \xmark\ ``A faster ball diffracts more.'' \cmark\ Faster means more momentum, a
        \emph{shorter} wavelength and \emph{less} diffraction.
  \item \xmark\ ``The evidence is that $\lambda$ equals the atom spacing.'' \cmark\ A similar
        size only lets diffraction happen. The evidence is that $\lambda$ found from the
        pattern matched $\dfrac{h}{mv}$.
  \item \xmark\ ``Davisson and Germer fired alpha particles at gold foil.'' \cmark\ That is
        Geiger--Marsden. Davisson and Germer fired electrons at a nickel crystal.
\end{itemize}
\end{mistakes}

\begin{examtip}
``Use de Broglie's idea to explain Bohr's stationary states'' (3 marks):
\begin{enumerate}
  \item Bohr's model gave no reason why only some orbits are allowed.
  \item De Broglie: a moving electron has a wavelength, $\lambda = \dfrac{h}{mv}$.
  \item In an allowed orbit a whole number of wavelengths fits ($2\pi r = n\lambda$): a
        standing wave, which does not radiate. In other orbits the wave cancels itself.
\end{enumerate}
``Explain how Davisson and Germer supported de Broglie'' (3 marks): a peak at one angle
$\to$ diffraction (only waves do this) $\to$ $\lambda$ from $d\sin\theta$ matched $h/mv$.
\end{examtip}

\begin{quickcheck}
  \item Find the de Broglie wavelength of a neutron ($\qty{1.675e-27}{kg}$) at
        \qty{2.2e3}{m.s^{-1}}. Could a crystal (atoms \qty{e-10}{m} apart) diffract it?
  \item An electron and a proton move at the same speed. Which has the longer wavelength, and
        how many times longer?
  \item A student says: ``A \qty{1200}{kg} car driving very slowly (\qty{1.0}{m.s^{-1}})
        through a garage door will spread out like a wave.'' Find its $\lambda$ and assess the
        claim.
  \item Electrons are accelerated from rest through \qty{150}{V}. Find their speed and
        wavelength. (Hint: $qV = \frac12 mv^2$, $q = \qty{1.602e-19}{C}$.)\label{ch3:qc:volt}
\end{quickcheck}

%% ==================================================================
\section{Bohr's limits and Schrödinger's quantum model}\label{ch3:sec:schrodinger}
\covers{M8.IQ3.01, M8.IQ3.05}

\begin{bigidea}
Bohr's model works only for hydrogen and misses finer details of spectra. In 1926 Erwin
Schrödinger treated the electron fully as a wave. His model gives the probability of finding
the electron, not an orbit, and it explains much more. It is the basis of today's model.
\end{bigidea}

\begin{keywords}
  \kw{Zeeman effect}{a spectral line splits when the atoms are in a magnetic field (an electric
    field gives the \emph{Stark effect}).}
  \kw{Fine structure}{one line that turns out to be two or more very close lines when you use
    a very good spectrometer.}
  \kw{Wave function ($\psi$)}{a maths description of the electron's wave. Where $\psi^2$ is
    big, the electron is likely to be found.}
  \kw{Orbital}{a 3-D region around the nucleus where the electron is likely to be found.}
  \kw{Quantum number}{a whole number that labels an allowed state of the electron, such as
    $n$.}
\end{keywords}

\subsection{What Bohr's model could not explain}

Bohr's model explained hydrogen's line spectrum and gave the Rydberg equation. But it has clear
limits (\cref{ch3:fig:zeeman}):
\begin{enumerate}
  \item It works only for hydrogen and for ions (charged atoms) that have just one electron,
        such as He\textsuperscript{+}.
  \item It cannot explain why some lines are brighter than others (their intensity).
  \item It cannot explain \term{fine structure}, the \term{Zeeman effect} or the Stark effect.
  \item It gives no reason for its rules, and it mixes classical orbits with quantum rules.
  \item It gives the electron an exact path, which the uncertainty principle (below) forbids.
\end{enumerate}

\begin{figure}
  \centering
  \begin{tikzpicture}[x=1cm, y=1cm, font=\small]
    \foreach \y/\lab in {1.6/{No field (normal view)}, 0.8/{In a magnetic field},
                         0/{No field, very high resolution}} {
      \fill[hscink] (0,\y) rectangle (5.5,\y+0.5);
      \node[anchor=east] at (-0.15,\y+0.25) {\lab};
    }
    \draw[red!70!white, line width=1.6pt] (2.75,1.6) -- (2.75,2.1);
    \foreach \x in {2.35, 2.75, 3.15} \draw[red!70!white, line width=1.4pt] (\x,0.8) -- (\x,1.3);
    \foreach \x in {2.68, 2.82} \draw[red!70!white, line width=1.1pt] (\x,0) -- (\x,0.5);
    \node[anchor=west, text=hscmuted] at (5.65,1.05) {Zeeman effect};
    \node[anchor=west, text=hscmuted] at (5.65,0.25) {fine structure};
  \end{tikzpicture}
  \caption{One spectral line seen three ways. Bohr's model predicts only the top picture.}
  \label{ch3:fig:zeeman}
\end{figure}

\subsection{What Schrödinger added}

\textbf{A wave equation.} Schrödinger took de Broglie's matter waves and wrote a
\term{wave equation} for the electron in 3-D. A wave equation says how a wave behaves. You
do not need the equation itself. Its solutions are \term{wave functions}, $\psi$ (``psi''). A
wave function is a maths description of the electron's wave. Max Born (1926) proposed that
$\psi^2$ gives the \term{probability} (how likely it is) of finding the electron at each
place. So we can only say where the electron is \emph{likely} to be.

\textbf{The uncertainty principle.} Heisenberg (1927) showed that you cannot know both the
exact position and the exact momentum of an electron at the same time. This is nature, not
poor equipment. So an electron has no definite orbit, which fits the probability picture.

\textbf{Energy levels come out of the maths.} A simple model shows the idea: an electron
trapped in a box (\cref{ch3:fig:box}). This model is not in the syllabus. It is a simple
picture of what Schrödinger's maths does. Only some waves fit, like on a guitar string:
\begin{itemize}
  \item the wave is zero at the walls, so only $\lambda = \dfrac{2L}{n}$ fits;
  \item $\lambda$ is $n$ times shorter, so $v$ is $n$ times bigger ($v = \dfrac{h}{m\lambda}$),
        and $K = \frac12 mv^2$ is $n^2$ times bigger: $E_1$, $4E_1$, $9E_1$, \ldots;
  \item for $n = 2$ the electron is never found in the exact middle.
\end{itemize}

\begin{figure}
  \centering
  \begin{tikzpicture}[x=1.1cm, y=0.62cm, font=\small]
    \fill[ground] (-0.25,-0.1) rectangle (0,4.3);
    \fill[ground] (4,-0.1) rectangle (4.25,4.3);
    \draw[hscink, line width=0.8pt] (0,-0.1) -- (0,4.3) (4,-0.1) -- (4,4.3);
    \foreach \n/\y/\lab in {1/0/{E_1}, 2/1.45/{4E_1}, 3/2.9/{9E_1}} {
      \draw[hscmuted, line width=0.4pt] (0,\y) -- (4,\y);
      \fill[hscseriesA!25] plot[domain=0:4, samples=120, smooth]
        (\x, {\y + 1.05*sin(\n*45*\x)^2}) -- (4,\y) -- (0,\y) -- cycle;
      \draw[hscseriesA, line width=1pt] plot[domain=0:4, samples=120, smooth]
        (\x, {\y + 1.05*sin(\n*45*\x)^2});
      \node[anchor=west] at (4.35,\y+0.35) {$n=\n$: $\lab$};
    }
    \draw[<->, hscink, line width=0.6pt] (0,-0.45) -- node[below] {box length $L$} (4,-0.45);
    \node[anchor=east, font=\footnotesize] at (-0.3,3.8) {wall};
    \node[anchor=west, font=\footnotesize] at (4.3,3.95) {wall};
  \end{tikzpicture}
  \caption{Electron in a box: the probability of finding it, for the three lowest energies.
    (The three states are drawn one above the other for clarity; the heights are not to
    energy scale.)}
  \label{ch3:fig:box}
\end{figure}

So quantised energy is a \emph{result} of the waves, not an assumed rule as in Bohr's model.
(In the box all the energy is kinetic. An atom is different: higher levels have slower
electrons and longer wavelengths, but more total energy, because the potential energy rises.)

\textbf{Orbitals, not orbits.} In an atom the standing waves are 3-D \term{orbitals}. We draw
them as electron clouds: the denser the dots, the more likely the electron is there
(\cref{ch3:fig:clouds}). Orbitals have names such as 1s and 2p. The number is the level $n$.
The letter gives the shape: s is a ball, and p is two blobs called \emph{lobes}.

Look at panel (b). The dots are most crowded near the nucleus. But think of the dots at one
distance as a thin shell, like the skin of a ball. A shell further out is much bigger than a
tiny region near the centre, so it holds more dots in total. If you count the dots at each
distance, the most likely distance is the \term{Bohr radius} (the dashed circle). It is the
radius of Bohr's $n = 1$ orbit, \qty{5.29e-11}{m}. So Schrödinger's model keeps Bohr's good
results.

\begin{figure}
  \centering
  \includegraphics{ch3_clouds.pdf}
  \caption{Bohr's fixed orbit and two of Schrödinger's orbitals. Each cloud shows where
    \emph{one} electron could be found. Panel (c) is drawn at a smaller scale.}
  \label{ch3:fig:clouds}
\end{figure}

\textbf{Quantum numbers come out of the maths.} Bohr's model has one number, $n$, and Bohr had
to assume it. Schrödinger's 3-D standing waves need three whole numbers, called
\term{quantum numbers}. They are $n$ (energy and size), $l$ (shape: s, p, d, \ldots) and $m_l$
(direction in space). All three come out of the wave equation; none is assumed. Later, two
more ideas were added. Electron \emph{spin} makes each electron act like a tiny magnet.
Pauli's \emph{exclusion principle} says that no two electrons in an atom can have the same set
of quantum numbers (spin included). Together they explain how electrons fill shells.

\textbf{What the new model explained.} It gives hydrogen's energy levels and the Rydberg
equation. It predicts how bright each line is. It explains the simplest splitting of lines in
a magnetic field (the Zeeman effect). It also works, with approximations, for atoms with many
electrons. Add electron spin and Pauli's exclusion principle, and it explains the periodic
table and chemical bonds. Add relativity (Paul Dirac, 1928), and it explains fine structure.
\Cref{ch3:tab:models} compares the models.

\begin{table}
  \caption{Three models of the atom.}\label{ch3:tab:models}
  \centering
  \small
  \begin{tabular}{@{}L{21mm}L{36mm}L{38mm}L{42mm}@{}}
    \toprule
    & \colhead{Rutherford (1911)} & \colhead{Bohr (1913)} & \colhead{Schrödinger (1926)}\\
    \midrule
    Electron & a particle & a particle & a wave, described by $\psi$\\
    Where & any orbit & fixed circular orbits & orbitals: probabilities only\\
    Energy & any value & quantised by a postulate & quantised by the wave maths\\
    Stable? & no: should collapse & yes, by assumption & yes: standing waves\\
    Quantum numbers & none & $n$ only (assumed) & $n$, $l$, $m_l$ (from the maths)\\
    Works for & $\alpha$ scattering only; no stable atom & hydrogen only & all atoms
      (approximately)\\
    \bottomrule
  \end{tabular}
\end{table}

\begin{workedexample}{An electron in a box (extension)}\label{ch3:we:box}
Extension (not examined): this simple model shows how wave maths gives quantised energy. An
electron is trapped in a box \qty{0.500}{nm} long. In its lowest state ($n = 1$), half a
wavelength fits in the box. Find its wavelength, speed and kinetic energy.
\given{$L = \qty{0.500}{nm} = \qty{5.00e-10}{m}$, \ $n = 1$, \ $m = \qty{9.109e-31}{kg}$,
  \ $h = \qty{6.626e-34}{J.s}$}
\find{$\lambda$, then $v$, then $K$}
\formula{$\lambda = \dfrac{2L}{n}$, \ $v = \dfrac{h}{m\lambda}$ (from $\lambda =
  \dfrac{h}{mv}$), \ $K = \frac12 mv^2$}
\substitute{\eqlines{\lambda &= \frac{2 \times \qty{5.00e-10}{m}}{1} = \qty{1.00e-9}{m}\\
  v &= \frac{\qty{6.626e-34}{J.s}}{\qty{9.109e-31}{kg} \times \qty{1.00e-9}{m}}
    = \qty{7.27e5}{m.s^{-1}}\\
  K &= \tfrac12 \times \qty{9.109e-31}{kg} \times (\qty{7.27e5}{m.s^{-1}})^2
    = \qty{2.41e-19}{J}\\
  K &= \frac{\qty{2.41e-19}{J}}{\num{1.602e-19}} = \qty{1.50}{eV} \quad\text{(J to eV)}}}
\answer{\final{$\lambda = \qty{1.00e-9}{m}$}, \ \final{$v = \qty{7.27e5}{m.s^{-1}}$}, \
  \final{$K = \qty{2.41e-19}{J} = \qty{1.50}{eV}$}}
\checkstep{a few eV is a typical atomic energy, so the model gives sensible sizes \cmark.
  The speed is far below $c$, so $\lambda = \dfrac{h}{mv}$ is fine.}
\end{workedexample}

\begin{workedexample}{Assessing Bohr's model (written answer)}\label{ch3:we:assess}
Plan and write a full-mark answer to this HSC-style question: ``Assess the limitations of
Bohr's model of the atom.'' (5 marks)
\given{the verb \emph{assess}: make a judgement of value, with reasons for and against}
\find{two successes, three limitations (each linked to evidence), and a judgement}
\formula{plan: successes $\to$ limitations $\to$ judgement}
\substitute{\textbf{Successes.} (1) It gives stable atoms, because it states (assumes) that
  electrons in stationary states do not radiate. (2) Its energy levels, with $hf = E_\mathrm{i} - E_\mathrm{f}$, give
  hydrogen's line wavelengths. For example, the jump $3 \to 2$ gives the measured red line at
  \qty{656}{nm}.\newline
  \textbf{Limitations.} (1) It works only for hydrogen and one-electron ions. It fails for the
  spectrum of helium. (2) It cannot explain why some lines are brighter than others, fine
  structure or the Zeeman effect. (3) Its rules are assumed, not explained, and it mixes
  classical orbits with quantum rules.}
\answer{\emph{Bohr's model was the first to put quantised energy into the atom, and it
  explained hydrogen's spectrum. But it failed for other atoms and for the finer details of
  spectra. So it was useful but incomplete, and Schrödinger's model later replaced it.}}
\checkstep{two successes $+$ three limitations linked to evidence $+$ a clear judgement
  \cmark. Each point is one short sentence.}
\end{workedexample}

\begin{mistakes}
\begin{itemize}
  \item \xmark\ ``Bohr explained why stationary states do not radiate.'' \cmark\ He
        \emph{assumed} it. De Broglie's standing waves later gave a reason.
  \item \xmark\ Calling fine structure ``the Stark effect''. \cmark\ Fine structure is seen with
        no field. Stark = electric field; Zeeman = magnetic field.
  \item \xmark\ ``The electron moves along a path inside the cloud'' or ``the cloud is many
        electrons''. \cmark\ There is no path; the cloud shows where \emph{one} electron is
        likely to be.
\end{itemize}
\end{mistakes}

\begin{examtip}
``Analyse Schrödinger's contribution to the current model of the atom'' (4--5 marks).
\emph{Analyse} means: name each part and say what it led to.
\begin{enumerate}
  \item A wave equation built on de Broglie's matter waves $\to$ energy levels and quantum
        numbers come out of the maths, not from postulates.
  \item The wave function $\psi$, with $\psi^2$ as probability (Born) $\to$ orbitals replace
        orbits. This fits the uncertainty principle.
  \item It works (approximately) for atoms with many electrons and predicts line brightness
        $\to$ it explains far more than Bohr's model.
  \item With spin, Pauli's principle and relativity (Dirac) added, it is the basis of today's
        quantum model, which is still being refined.
\end{enumerate}
The 2025 HSC exam asked for \emph{two} ways in which Schrödinger's model of electron
behaviour differs from the Rutherford and Bohr models (3 marks). Contrast each point, for
example: \emph{electrons are waves described by probabilities, not particles on fixed orbits};
\emph{quantised energy comes from the wave equation, not from a postulate}.
\end{examtip}

\begin{quickcheck}
  \item State three observations that Bohr's model cannot explain.
  \item What is the difference between the Zeeman effect and the Stark effect?
  \item Outline two ways in which Schrödinger's model of the atom differs from Bohr's.
  \item (Extension) For the electron in \cref{ch3:we:box}: (a) find the kinetic energy of the
        $n = 2$ state; (b) the electron drops from $n = 2$ to $n = 1$. Find the energy (in eV)
        and the wavelength of the photon it emits. In which part of the EM spectrum is it?
\end{quickcheck}

%% ==================================================================
\section{Key experiments}\label{ch3:sec:experiments}
\covers{M8.IQ3.02, M8.IQ3.04}

\begin{refresher}{Diffraction gratings (Module 7)}
A \term{diffraction grating} has thousands of narrow slits, a distance $d$ apart. At some
angles, the waves from neighbouring slits arrive in step and add up. This is
\term{constructive interference}, and we see a bright line. Bright lines obey
$d\sin\theta = m\lambda$ (on the formula sheet). Here $m$ is the order of the line (a whole
number), not a mass. $m = 0$ is the straight-through image (the \emph{zero order}). $m = 1$ is the first bright line on each side (the \emph{first order}).
For the first order, $\lambda = d\sin\theta$. A crystal's rows of atoms act in the same way.
\end{refresher}

\begin{toolkit}{Why the gradient of the Balmer graph is $R$}
A straight line through the origin has the form $y = mx$, where $m$ is the gradient. For the
Balmer series ($n_\mathrm{f} = 2$) the Rydberg equation has the same form:
\[ \underbrace{\frac{1}{\lambda}}_{y} = \underbrace{R}_{m} \times
   \underbrace{\left(\frac{1}{2^2} - \frac{1}{n_\mathrm{i}^2}\right)}_{x} \]
So the gradient of this graph is $R$. To find it, pick two points far apart \emph{on the line
of best fit} (not data points). Then gradient $= \dfrac{\text{rise}}{\text{run}} =
\dfrac{y_2 - y_1}{x_2 - x_1}$. In \cref{ch3:fig:balmergraph}, the points $(0,\ 0)$ and $(0.25,\ \qty{2.74e6}{m^{-1}})$ give
\[ \text{gradient} = \frac{\qty{2.74e6}{m^{-1}} - 0}{0.25 - 0} = \qty{1.10e7}{m^{-1}} \]
$x$ has no unit, so the gradient has the unit of $y$ (\si{m^{-1}}).
\end{toolkit}

\begin{experiment}{Investigating the Balmer series of hydrogen}
\aim{To measure the wavelengths of hydrogen's visible lines and compare them with the Rydberg
  equation for jumps to $n = 2$.}
\apparatus{Hydrogen \term{discharge tube} (a glass tube of low-pressure hydrogen that glows
  when a high voltage is applied) and its supply (about \qty{5}{kV}). A spectrometer: a
  collimator (a tube with a narrow slit that makes a thin, parallel beam), a turning telescope
  and an angle scale. A diffraction grating with 600 lines per mm, so
  $d = \dfrac{\qty{1}{mm}}{600} = \qty{1.667e-6}{m}$. A darkened room.
  \begin{figure}
    \centering
    \begin{tikzpicture}[x=0.9cm, y=0.9cm, font=\small]
      \fill[hscfield!25] (-5.2,0) circle (0.32);
      \draw[hscfield, line width=0.8pt] (-5.2,0) circle (0.32);
      \node[align=center, below] at (-5.2,-0.4) {hydrogen\\tube};
      \draw[hscink, line width=0.8pt] (-4.4,-0.3) rectangle (-1.6,0.3);
      \draw[hscink, line width=2pt] (-4.4,-0.3) -- (-4.4,-0.06) (-4.4,0.06) -- (-4.4,0.3);
      \node[above] at (-4.4,0.3) {slit};
      \node[below] at (-3.0,-0.3) {collimator};
      \draw[hscray, line width=1pt, ->] (-1.5,0) -- (-0.12,0);
      \draw[hscmuted, dashed] (0,0) circle (1.1);
      \draw[hscink, line width=1.6pt] (0,-0.7) -- (0,0.7);
      \node[align=center, anchor=north] at (0,-1.15) {grating};
      \draw[guide] (0,0) -- (4.6,0);
      \node[anchor=west] at (4.65,0) {zero order};
      \begin{scope}[rotate=23.2]
        \draw[hscray, line width=1pt, ->] (0.12,0) -- (1.9,0);
        \draw[hscink, line width=0.8pt, fill=white] (2.0,-0.28) rectangle (4.5,0.28);
        \node[rotate=23.2] at (3.25,0) {telescope};
      \end{scope}
      \begin{scope}[rotate=-23.2]
        \draw[hscink, dashed, line width=0.6pt] (2.0,-0.28) rectangle (4.5,0.28);
        \node[text=hscmuted, anchor=west] at (4.55,0) {other side};
      \end{scope}
      \draw[->, hscink, line width=0.6pt] (1.5,0) arc[start angle=0, end angle=23.2, radius=1.5];
      \node at (10:1.8) {$\theta$};
    \end{tikzpicture}
    \caption{Top view of a grating spectrometer. Measure $\theta$ on both sides.}
    \label{ch3:fig:balmersetup}
  \end{figure}}
\method{\begin{enumerate}
  \item In a dark room, put the glowing tube by the slit. Set the grating at right angles to
        the light.
  \item Turn the telescope to the first-order red line and record the angle. Do the same on
        the other side; $\theta$ is half the angle between the two positions.
  \item Repeat for the cyan, blue-violet and violet lines, three times each.
  \item Calculate $\lambda = d\sin\theta$ ($m = 1$). Plot $\frac{1}{\lambda}$ against
        $\left(\frac{1}{2^2} - \frac{1}{n_\mathrm{i}^2}\right)$.
\end{enumerate}
\textbf{Safety:} high voltage: do not touch the tube's ends. The tube gets hot and gives some
UV, so do not stare at it.}
\results{Sample results (angles $\pm\ang{0.1}$):
  \begin{center}\small
  \begin{tabular}{@{}lccccc@{}}
    \toprule
    Line & $n_\mathrm{i}$ & $\theta$ & $\lambda = d\sin\theta$ (nm) & Rydberg $\lambda$ (nm)
      & Difference\\
    \midrule
    red         & 3 & \ang{23.3} & 659.2 & 656.3 & $+0.4\%$\\
    cyan        & 4 & \ang{17.0} & 487.3 & 486.2 & $+0.2\%$\\
    blue-violet & 5 & \ang{15.0} & 431.4 & 434.1 & $-0.6\%$\\
    violet      & 6 & \ang{14.3} & 411.7 & 410.2 & $+0.4\%$\\
    \bottomrule
  \end{tabular}
  \end{center}
  Difference $= \dfrac{\text{measured} - \text{Rydberg}}{\text{Rydberg}} \times 100\%$. For
  example, red: $\dfrac{659.2 - 656.3}{656.3} \times 100\% = +0.4\%$.
  \begin{figure}
  \centering
  \begin{tikzpicture}
    \begin{axis}[hsc, width=0.58\linewidth, height=0.32\linewidth,
        xmin=0, xmax=0.25, ymin=0, ymax=2.8,
        xlabel={$\frac{1}{2^2} - \frac{1}{n_\mathrm{i}^2}$},
        ylabel={$\frac{1}{\lambda}$ ($10^{6}$ \si{m^{-1}})},
        xtick={0,0.05,0.10,0.15,0.20,0.25},
        xticklabel style={/pgf/number format/.cd, fixed, precision=2},
        ytick={0,0.5,...,2.5}, yticklabel style={/pgf/number format/.cd, fixed, precision=1}]
      \addplot[hscseriesA, only marks, mark=*, mark size=2pt]
        coordinates {(0.13889,1.5170) (0.1875,2.0521) (0.21,2.3180) (0.22222,2.4290)};
      \addplot[hscseriesB, dashed, line width=1pt, domain=0:0.25, samples=2] {10.965*x};
      \node[anchor=west, font=\small, text=hscink, align=left] at (axis cs:0.01,2.45)
        {line through the origin:\\gradient $= \qty{1.10e7}{m^{-1}}$};
    \end{axis}
  \end{tikzpicture}
  \caption{Sample Balmer data. The gradient is the Rydberg constant (see the Maths toolkit
    before this experiment).}\label{ch3:fig:balmergraph}
  \end{figure}}
\conclusion{All four wavelengths agree with the Rydberg values to within 1\%. The graph is a
  straight line through the origin with gradient \qty{1.10e7}{m^{-1}}, close to $R =
  \qty{1.097e7}{m^{-1}}$. So each visible line is a jump down to $n = 2$ from $n = 3$, 4, 5 or
  6. This supports quantised energy levels.}
\significance{These are the lines Bohr's model was built to explain. Astronomers use the same
  lines to find hydrogen in stars and gas clouds.}
\validity{\begin{itemize}
  \item \textbf{Validity:} use a pure hydrogen tube and ignore faint impurity lines. Keep the
        grating at right angles to the beam. Use first-order lines and the same grating.
  \item \textbf{Reliability:} repeat each angle three times and average; the two sides and
        other groups should agree.
  \item \textbf{Accuracy:} compare with accepted values (percentage difference). Reading angles
        to \ang{0.1} gives $\lambda$ to about \qty{3}{nm}. Measuring both sides cancels a zero
        error (a scale that does not read exactly zero when it should). A hand spectroscope is
        less accurate.
\end{itemize}}
\end{experiment}

\begin{experiment}{Davisson and Germer: electrons diffract (1927)}
\aim{To study how electrons bounce off a nickel crystal. The results became the first direct
  test of de Broglie's matter waves.}
\apparatus{Electron gun (hot filament, \qty{54}{V} accelerating voltage for the key result),
  vacuum chamber, nickel crystal, and a detector that moves on an arc and measures the electron
  current at each angle.
  \begin{figure}
    \centering
    \begin{tikzpicture}[x=0.9cm, y=0.9cm, font=\small]
      \draw[hscmuted, dashed, rounded corners=6pt] (-3.4,-1.35) rectangle (3.8,4.3);
      \node[anchor=north east, text=hscmuted] at (3.75,4.25) {vacuum};
      \fill[hscmuted!35] (-1.6,-0.6) rectangle (1.6,0);
      \foreach \x in {-1.4,-1.0,...,1.45} \fill[hscink] (\x,-0.15) circle (1.3pt);
      \node[anchor=north] at (0,-0.6) {nickel crystal};
      \draw[hscink, line width=0.8pt, fill=white] (-0.35,3.35) rectangle (0.35,4.0);
      \draw[hscink, decorate, decoration={zigzag, amplitude=1pt, segment length=3pt}]
        (-0.2,3.8) -- (0.2,3.8);
      \node[anchor=east, align=right] at (-0.45,3.65) {electron gun\\(\qty{54}{V})};
      \draw[hscvel, line width=1.2pt, ->] (0,3.3) -- (0,0.05);
      \draw[guide] (0,0) -- (0,2.9);
      \draw[hscvel, line width=1.2pt, ->] (0,0) -- (40:2.45);
      \draw[hscmuted, dashed] (78:2.85) arc[start angle=78, end angle=25, radius=2.85];
      \begin{scope}[shift={(40:2.85)}, rotate=40]
        \draw[hscink, line width=0.8pt, fill=hscseriesB!40] (-0.25,-0.3) rectangle (0.35,0.3);
      \end{scope}
      \node[anchor=west, align=left] at (2.75,0.95) {movable\\detector};
      \draw[hscink, line width=0.6pt] (0,1.0) arc[start angle=90, end angle=40, radius=1.0];
      \node at (65:1.3) {$\theta$};
    \end{tikzpicture}
    \hspace{4mm}
    \begin{tikzpicture}
      \begin{axis}[hsc, width=0.46\linewidth, height=0.36\linewidth,
          xmin=0, xmax=90, ymin=0, ymax=2.0, xtick={0,30,60,90}, ytick=\empty,
          xlabel={angle $\theta$ (degrees)}, ylabel={electrons detected}]
        \addplot[hscseriesA, line width=1.2pt, domain=0:90, samples=181, smooth]
          {0.55 + 0.004*x + 1.0*exp(-((x-50)/6.5)^2)};
        \node[anchor=west, font=\small] at (axis cs:55,1.6) {peak at \ang{50}};
      \end{axis}
    \end{tikzpicture}
    \caption{The Davisson--Germer apparatus (the beam hits the crystal head-on) and a sketch of
      the result for \qty{54}{V} electrons.}
    \label{ch3:fig:dg}
  \end{figure}}
\method{Fire electrons of known energy (fixed voltage) straight at the crystal. Move the
  detector step by step and record the current at each angle. Repeat for several voltages.
  (By accident, air got in and spoiled the nickel surface. To clean it, they heated it and
  then let it cool slowly. This is called \term{annealing}. It turned the nickel into a few
  large, regular crystals. Their rows of atoms act like a diffraction grating.)}
\results{A strong peak at $\theta = \ang{50}$ for \qty{54}{V} electrons (\cref{ch3:fig:dg}).
  The peak moved when the voltage changed. Compare two wavelengths. The row spacing,
  $d = \qty{2.15e-10}{m}$, was measured separately by X-ray diffraction.
  \begin{itemize}
    \item From the diffraction: $\lambda = d\sin\theta = \qty{2.15e-10}{m} \times
          \sin\ang{50} = \qty{1.65e-10}{m}$.
    \item From de Broglie: kinetic energy $qV = \qty{1.602e-19}{C} \times \qty{54}{V} =
          \qty{8.65e-18}{J}$. From $qV = \frac12 mv^2$, $v = \sqrt{2qV/m}$ (not on the sheet):
          \begin{align*}
            v &= \sqrt{\frac{2 \times \qty{8.65e-18}{J}}{\qty{9.109e-31}{kg}}}
              = \qty{4.36e6}{m.s^{-1}}\\
            \lambda &= \frac{h}{mv} = \frac{\qty{6.626e-34}{J.s}}{\qty{9.109e-31}{kg} \times
              \qty{4.36e6}{m.s^{-1}}} = \qty{1.67e-10}{m}
          \end{align*}
  \end{itemize}}
\conclusion{A peak at one angle is constructive interference, so the electrons were diffracted:
  they behaved as waves. The two wavelengths agree to about 1\%, which confirms
  $\lambda = \dfrac{h}{mv}$.}
\significance{The first direct evidence of matter waves. Nobel Prizes: de Broglie (1929);
  Davisson and G.\,P. Thomson (1937). Electron microscopes use the same idea.}
\validity{\begin{itemize}
  \item \textbf{Validity:} the vacuum stops collisions with air; a fixed voltage gives a known
        electron energy; a single crystal gives regular rows of atoms. The peak moved with
        voltage as de Broglie predicted.
  \item \textbf{Reliability:} the angles were swept many times and the peak appeared each time.
  \item \textbf{Accuracy:} the two wavelengths agree to about 1\%; the main uncertainties are
        in $\theta$ and $d$.
\end{itemize}}
\end{experiment}

\begin{experiment}{G.\,P. Thomson: electrons through thin films (1927)}
\aim{To test whether fast electrons are diffracted when they pass through a very thin film.}
\apparatus{Electron gun (tens of kilovolts), very thin film (celluloid at first, then metals
  such as gold and aluminium), vacuum tube, photographic plate. School version: an
  electron-diffraction tube with a thin graphite film and a \qtyrange{2}{5}{kV} supply.
  \begin{figure}
    \centering
    \begin{tikzpicture}[x=0.85cm, y=0.75cm, font=\small]
      \draw[hscink, line width=0.8pt, fill=white] (-4.2,-0.25) rectangle (-3.0,0.25);
      \node[below] at (-3.6,-0.35) {electron gun};
      \draw[hscvel, line width=1.2pt, ->] (-3.0,0) -- (-0.1,0);
      \fill[hscmuted] (-0.05,-0.8) rectangle (0.05,0.8);
      \node[above] at (0,0.85) {thin film};
      \foreach \a in {9,16} {
        \draw[hscvel, line width=0.7pt] (0,0) -- (3.2,{3.2*tan(\a)});
        \draw[hscvel, line width=0.7pt] (0,0) -- (3.2,{-3.2*tan(\a)});
      }
      \draw[hscvel, line width=0.7pt] (0,0) -- (3.2,0);
      \draw[hscink, line width=1.2pt] (3.2,-1.2) -- (3.2,1.2);
      \node[below] at (3.2,-1.25) {screen};
      \begin{scope}[shift={(6.0,0)}, x=0.75cm]
        \fill[hscink] (0,0) circle (1.2);
        \foreach \r in {0.5,0.92} \draw[hscseriesC!70!white, line width=1.4pt] (0,0) circle (\r);
        \fill[hscseriesC!70!white] (0,0) circle (2pt);
        \node[below, align=center] at (0,-1.25) {rings seen\\on the screen};
      \end{scope}
    \end{tikzpicture}
    \caption{Electrons passing through a thin film make rings on the screen.}
    \label{ch3:fig:thomson}
  \end{figure}}
\method{Fire the beam through the film and record the pattern. Change the voltage and watch the
  rings. Bring a magnet near: if the pattern moves, it is made by charged electrons, not X-rays.}
\results{Bright rings around a central spot, like the rings X-rays (known waves) make. The
  rings got smaller at higher voltage, and the magnet moved the pattern. They are rings, not
  spots, because the film is many tiny crystals facing all directions. So each diffracted beam
  spreads into a cone, which hits the screen as a ring.}
\conclusion{Rings are a diffraction pattern, so electrons behave as waves. A higher voltage
  gives faster electrons, a shorter wavelength ($\lambda = h/mv$), less spreading and smaller
  rings. The ring sizes matched de Broglie's wavelength.}
\significance{It confirmed matter waves by a different method, and schools can repeat it.}
\validity{\begin{itemize}
  \item \textbf{Validity:} use a vacuum and a very thin film; the magnet test shows the
        pattern is from electrons.
  \item \textbf{Reliability:} measure each ring several times at each voltage; average.
  \item \textbf{Accuracy:} the rings are wide and fuzzy, so use several voltages and look for
        the trend.
\end{itemize}}
\end{experiment}
